%!TEX root = ../QuantumNoiseAndDecoherence.tex
% ──────────────────────────────────────────────────────────────
%  Lecture 6 — Gaussian Models and Symplectic Diagonalisation
% ──────────────────────────────────────────────────────────────

% Continues Section 4 (Continuous Variables).

\subsection{Quadratic Hamiltonians and Gaussian models}%
\label{sec:gaussian-models}

To model the environment~$E$ in our system--environment picture, we
need exactly solvable models for bosonic fields.  \emph{Quadratic}
(Gaussian) Hamiltonians fill this role: they are general enough to
capture physics from phonon propagation in solids to relativistic
quantum field theory, yet they can be diagonalised using only linear
algebra.

\begin{definition}[Gaussian (quasi-free) Hamiltonian]%
\label{def:gaussian-hamiltonian}
  Given $n$~modes with canonical operators
  $\hat{x}_1,\ldots,\hat{x}_n$, $\hat{p}_1,\ldots,\hat{p}_n$
  satisfying the CCR
  $[\hat{x}_j,\,\hat{p}_k] = i\delta_{jk}\One$, the
  \textbf{Gaussian Hamiltonian} is
  \begin{equation}\label{eq:gaussian-H}
    \eqbox{\hat{H}
    = \tfrac{1}{2}\sum_{j,k=1}^{n}
      \bigl[V_x\bigr]_{jk}\,\hat{x}_j\,\hat{x}_k
    + \tfrac{1}{2}\sum_{j,k=1}^{n}
      \bigl[V_p\bigr]_{jk}\,\hat{p}_j\,\hat{p}_k}\,,
  \end{equation}
  where $V_x, V_p \in \R^{n\times n}$ are symmetric, positive-definite
  matrices encoding the potential and kinetic energy, respectively.
\end{definition}

\begin{intuition}[Scope of the Gaussian model]
  Despite having only quadratic terms, this Hamiltonian covers:
  $V_x$, $V_p$ diagonal $\Rightarrow$ independent oscillators;
  $V_x$ off-diagonal, $V_p = \One$ $\Rightarrow$ phonons in solids;
  both off-diagonal $\Rightarrow$ the Klein--Gordon field on a lattice.
  The mode number~$n$ can be arbitrarily large; all results below
  remain valid.
\end{intuition}

\subsubsection{Compact notation}

Introduce the column vector of $2n$~canonical operators
\begin{equation}\label{eq:r-vector}
  \hat{\vb{r}}
  = (\hat{x}_1,\ldots,\hat{x}_n,\,
     \hat{p}_1,\ldots,\hat{p}_n)^\top
\end{equation}
and the $2n \times 2n$ Hamiltonian matrix
\begin{equation}\label{eq:H-matrix}
  h = \begin{pmatrix} V_x & 0 \\ 0 & V_p \end{pmatrix}\,.
\end{equation}
Then the Hamiltonian takes the compact form
\begin{equation}\label{eq:H-compact}
  \hat{H} = \tfrac{1}{2}\,\hat{\vb{r}}^\top h\,\hat{\vb{r}}\,.
\end{equation}

In terms of ladder operators
$\hat{b}_j = (\hat{x}_j + i\hat{p}_j)/\sqrt{2}$, the Hamiltonian
contains number-conserving terms ($\hat{b}_j^\adj \hat{b}_k$) and
\emph{squeezing} terms ($\hat{b}_j\hat{b}_k$,
$\hat{b}_j^\adj\hat{b}_k^\adj$) that do not conserve particle number.

% ──────────────────────────────────────────────────────────────
\subsection{Symplectic transformations}%
\label{sec:symplectic}

The strategy for diagonalising~\eqref{eq:H-compact} is to make a
linear change of variables on the operators that preserves the CCR.

\begin{definition}[Symplectic form and symplectic group]%
\label{def:symplectic}
  The $2n \times 2n$ \textbf{symplectic form} is
  \begin{equation}\label{eq:sigma}
    \sigma = \begin{pmatrix} 0 & \One_n \\ -\One_n & 0 \end{pmatrix}\,.
  \end{equation}
  The canonical commutation relations in matrix form read
  \begin{equation}\label{eq:ccr-matrix}
    [\hat{r}_j,\, \hat{r}_k] = i\,\sigma_{jk}\,\One\,.
  \end{equation}
  A real $2n \times 2n$ matrix $S$ is \textbf{symplectic} if it
  preserves the symplectic form:
  \begin{equation}\label{eq:symplectic-cond}
    \eqbox{S\,\sigma\, S^\top = \sigma}\,.
  \end{equation}
  The set of all such matrices forms the \textbf{real symplectic group}
  $\mathrm{Sp}(2n,\R)$.
\end{definition}

Defining $\hat{\vb{r}}' = S\,\hat{\vb{r}}$ with
$S \in \mathrm{Sp}(2n,\R)$ preserves the CCR.  Our task reduces to
finding an~$S$ that brings~$h$ to diagonal form.

% ──────────────────────────────────────────────────────────────
\subsection{Williamson's theorem}%
\label{sec:williamson}

\begin{theorem}[Williamson's normal form]\label{thm:williamson}
  Let $M$ be a $2n \times 2n$ real, symmetric, positive-definite
  matrix.  Then there exists a symplectic matrix
  $S \in \mathrm{Sp}(2n,\R)$ such that
  \begin{equation}\label{eq:williamson}
    \eqbox{S\, M\, S^\top
    = \diag(\mu_1,\ldots,\mu_n,\,\mu_1,\ldots,\mu_n)}\,,
  \end{equation}
  where the \textbf{symplectic eigenvalues}
  $\mu_1 \geq \mu_2 \geq \cdots \geq \mu_n > 0$ are uniquely
  determined by~$M$.
\end{theorem}

\begin{remark}
  The symplectic eigenvalues can be found without constructing~$S$:
  they are the positive eigenvalues of the matrix $i\sigma M$, or
  equivalently the positive square roots of the eigenvalues of
  $-(\sigma M)^2$.
\end{remark}

\subsubsection{Constructing $S$ for the Gaussian Hamiltonian}

For $h = \diag(V_x, V_p)$, an explicit symplectic diagonaliser can be
constructed.  Define the $n \times n$ matrix
\begin{equation}\label{eq:G-matrix}
  G = V_p^{1/2}\, V_x\, V_p^{1/2}\,,
\end{equation}
which is symmetric and positive definite.  Let $Q$ be an orthogonal
matrix diagonalising~$G$:
\begin{equation}\label{eq:Q-diag}
  Q^\top G\, Q = D^2
  = \diag(\mu_1^2,\ldots,\mu_n^2)\,,
\end{equation}
where $\mu_j > 0$ are the symplectic eigenvalues of~$h$.

\begin{keyresult}[Explicit symplectic diagonaliser]
  The matrix
  \begin{equation}\label{eq:S-explicit}
    S = \begin{pmatrix}
      D^{1/2}\, Q^\top V_x^{-1/2} & 0 \\[4pt]
      0 & D^{-1/2}\, Q^\top V_p^{-1/2}
    \end{pmatrix}
  \end{equation}
  is symplectic and satisfies $S\, h\, S^\top = \diag(\mu_1,\ldots,\mu_n,\,\mu_1,\ldots,\mu_n)$.
\end{keyresult}

\begin{exercise}\label{ex:verify-S}
  Verify that~\eqref{eq:S-explicit} satisfies $S\sigma S^\top = \sigma$
  and $ShS^\top = \diag(\mu_1,\ldots,\mu_n,\mu_1,\ldots,\mu_n)$.
\end{exercise}

\begin{intuition}[When is this easy?]
  In non-relativistic physics $V_p = \One$, so $G = V_x$.  For
  translation-invariant systems $V_x$ is circulant and diagonalisable
  by the discrete Fourier transform, making the construction analytic.
\end{intuition}

% ──────────────────────────────────────────────────────────────
\subsection{Diagonalised form and normal modes}%
\label{sec:normal-modes}

Defining new canonical variables
$\hat{\vb{r}}' = S\,\hat{\vb{r}}$ (which still obey the CCR), the
Hamiltonian becomes
\begin{equation}\label{eq:H-diagonal}
  \hat{H}
  = \sum_{j=1}^{n} \frac{\mu_j}{2}
    \bigl(\hat{x}_j'^{\,2} + \hat{p}_j'^{\,2}\bigr)
  = \sum_{j=1}^{n} \mu_j
    \Bigl(\hat{a}_j^\adj \hat{a}_j + \tfrac{1}{2}\Bigr)\,,
\end{equation}
where $\hat{a}_j = (\hat{x}'_j + i\hat{p}'_j)/\sqrt{2}$ are the
\textbf{normal-mode} annihilation operators.  This is a collection of
$n$~independent harmonic oscillators with frequencies
$\mu_1,\ldots,\mu_n$.

\begin{keyresult}[Reduction to independent oscillators]
  Every Gaussian Hamiltonian~\eqref{eq:gaussian-H} is unitarily
  equivalent to a sum of independent harmonic oscillators---justifying
  the name \emph{quasi-free}.
\end{keyresult}

% ──────────────────────────────────────────────────────────────
\subsection{Transformation of the characteristic function}%
\label{sec:char-fn-transform}

The symplectic change of variables acts simply on the characteristic
function: if $\chi_\rho(\vb{\xi}) = \tr(\rho\, D(\vb{\xi}))$, then
\begin{equation}\label{eq:chi-transform}
  \chi_\rho(\vb{\xi})
  = \chi_{\rho'}\!\bigl((S^{-1})^\top \vb{\xi}\bigr)\,,
\end{equation}
where $\rho'$ is the state in the normal-mode basis.
This implies simple transformation rules for the moments:
\begin{equation}\label{eq:moments-transform}
  \eqbox{\vb{m} = S^{-1}\,\vb{m}'\,,
  \qquad
  \gamma = S^{-1}\,\gamma'\,(S^{-1})^\top}\,,
\end{equation}
where $\vb{m} = \expect{\hat{\vb{r}}}$ and
$\gamma_{jk} = \tfrac{1}{2}\expect{\{\Delta\hat{r}_j,\,\Delta\hat{r}_k\}}$
with $\Delta\hat{r}_j = \hat{r}_j - m_j$.

% ──────────────────────────────────────────────────────────────
\subsection{The thermal state of a Gaussian model}%
\label{sec:thermal-state}

Since $\hat{H}$ is a sum of independent oscillators in the
normal-mode basis, the thermal state
$\rho_\beta = e^{-\beta \hat{H}}/Z$ factorises:
\begin{equation}\label{eq:thermal-factorised}
  \rho_\beta
  = \bigotimes_{j=1}^{n} \frac{e^{-\beta\mu_j(\hat{a}_j^\adj \hat{a}_j + 1/2)}}{Z_j}\,,
  \qquad
  Z_j = \frac{1}{2\sinh(\beta\mu_j/2)}\,.
\end{equation}
Each factor is Gaussian, so the full thermal state is Gaussian with
first moments $\vb{m} = 0$ (by symmetry) and covariance matrix:

\begin{keyresult}[Covariance matrix of the thermal state]
  In the normal-mode variables, the covariance matrix is diagonal:
  \begin{equation}\label{eq:gamma-prime}
    \gamma'
    = \diag\!\bigl(\bar{n}_1 + \tfrac{1}{2},\ldots,\bar{n}_n + \tfrac{1}{2},\,
      \bar{n}_1 + \tfrac{1}{2},\ldots,\bar{n}_n + \tfrac{1}{2}\bigr)\,,
  \end{equation}
  where $\bar{n}_j = (e^{\beta\mu_j} - 1)^{-1}$ is the Bose--Einstein
  mean occupation number.  In the original variables, the covariance
  matrix is
  \begin{equation}\label{eq:gamma-beta}
    \eqbox{\gamma_\beta = S^{-1}\,\gamma'\,(S^{-1})^\top}\,.
  \end{equation}
\end{keyresult}

\begin{remark}
  Proving that a single-mode thermal state has a Gaussian
  characteristic function requires summing Laguerre polynomials
  over Fock states and applying a generating-function identity.
  The result:
  $\chi(\xi) \propto \exp\!\bigl(-(\bar{n}+\tfrac{1}{2})\lvert\xi\rvert^2\bigr)$.
\end{remark}

\begin{exercise}\label{ex:thermal-covariance}
  For $n$~independent oscillators with $V_x = \diag(\omega_1^2,\ldots,\omega_n^2)$
  and $V_p = \One$, show that:
  \begin{enumerate}[(a)]
    \item The symplectic eigenvalues are $\mu_j = \omega_j$.
    \item The covariance matrix in the original variables is
      $\gamma_\beta = \diag\bigl((\bar{n}_1+\frac{1}{2})/\omega_1,\ldots,
      (\bar{n}_n+\frac{1}{2})/\omega_n,\,
      (\bar{n}_1+\frac{1}{2})\omega_1,\ldots,
      (\bar{n}_n+\frac{1}{2})\omega_n\bigr)$.
  \end{enumerate}
\end{exercise}

\begin{intuition}[Why this result matters]
  The covariance matrix~\eqref{eq:gamma-beta} is the starting point
  for modelling the environment.  Combined with Wick's theorem it
  determines all correlation functions of the thermal field, and its
  entries will later set the decay rates in the Lindblad equation.
\end{intuition}
